\begin{definition}[Known: conventions]\label{A:conventions}
For a small category $A$ we put $BA=|NA|$. For $H:A\to\Set$, the category $\El_AH$ has objects $(a,h)$ and arrows $f:(a,h)\to(b,k)$ satisfying $H(f)h=k$. We call $A$ total when it is nonempty and connected and every component of every $\El_AH$ is weakly contractible. We call $A$ exact when $\colim_A:\Ab^A\to\Ab$ is exact. We write $P_x=\Z[A(x,-)]$ and $T_i(M)=L_i\colim_A M$.
\end{definition}
\begin{unresolved}[Main question]\label{A:main}
Does every total small category have exact abelian colimits? No proof or counterexample is asserted at this checkpoint.
\end{unresolved}

\begin{theorem}[Proved: detection by all prime fields]\label{A:fields}
Let $A$ be any small category, possibly empty or disconnected. The following are equivalent:
\begin{enumerate}[label=(\roman*)]
\item $\colim_A:\Ab^A\to\Ab$ is exact.
\item Colimits over $A$ are exact for $\Q$-vector-space diagrams and for $\mathbb F_p$-vector-space diagrams for every prime $p$.
\end{enumerate}
Consequently the conjunction of all coefficient implications requested in Part A.2 is equivalent to Part A.1. This is a statement about all primes together; it makes no assertion that finitely many field tests suffice.
\end{theorem}
\begin{proof}
The underlying abelian group of a colimit of $R$-module diagrams is its abelian-group colimit, with its uniquely induced $R$-action. Indeed, the usual direct-sum presentation and its relation subgroup are $R$-linear. Thus (i) implies (ii).

Assume (ii). For a covariant abelian diagram $M$, let $C_\bullet(A,M)$ be its bar complex, with
\[
C_n(A,M)=\bigoplus_{a_0\to\cdots\to a_n}M(a_0).
\]
The first face applies the transition map of $M$, the intermediate faces compose arrows, and the last face omits the last arrow; degeneracies insert identities. Its homology is $T_n(M)$. This statement applies even when the groups $M(a)$ are not projective: an objectwise projective resolution in $\Ab^A$, followed by the bar construction, gives a first-quadrant double complex; direct sums are exact in the resolution direction, and the augmented bar complex contracts on each representable projective in the other direction. The two filtrations give the asserted comparison.

First suppose every $M(a)$ is torsionfree. The complex $C_\bullet=C_\bullet(A,M)$ consists of flat $\Z$-modules, and tensoring it with a field $k$ gives literally $C_\bullet(A,M\otimes_\Z k)$. By (ii), its positive homology vanishes for $k=\Q$ and for every $k=\mathbb F_p$. Flat base change to $\Q$ gives
\[
H_i(C_\bullet)\otimes\Q=0\qquad(i>0),
\]
so each $H_i(C_\bullet)$ is torsion. For a complex of flat $\Z$-modules the universal-coefficient sequence is
\[
0\longrightarrow H_j(C_\bullet)\otimes\mathbb F_p
\longrightarrow H_j(C_\bullet\otimes\mathbb F_p)
\longrightarrow \Tor^{\Z}_1(H_{j-1}(C_\bullet),\mathbb F_p)
\longrightarrow0.
\]
For completeness, this follows by tensoring the cycle and boundary short exact sequences; their boundary groups are torsionfree subgroups of flat abelian groups and hence flat. Taking $j=i+1$ shows that $H_i(C_\bullet)[p]=0$ for every prime $p$. A nonzero torsion abelian group contains an element of prime order, so $H_i(C_\bullet)=0$. Thus every objectwise torsionfree diagram is colimit-acyclic.

Next suppose $M$ is objectwise torsion. It decomposes naturally as the direct sum of its $p$-primary subfunctors. A $p$-primary diagram is the filtered union of the subfunctors $M[p^n]$. Each quotient $M[p^n]/M[p^{n-1}]$ is killed by $p$, and therefore is a diagram of $\mathbb F_p$-vector spaces. Its integral bar complex is the underlying abelian-group complex of its $\mathbb F_p$ bar complex, so (ii) makes it acyclic. Induction on $n$ using the long exact sequence of derived colimits proves acyclicity of $M[p^n]$. The bar complex and homology commute with direct sums and filtered colimits of abelian groups, proving acyclicity of $M$.

For arbitrary $M$, its torsion subgroups form a subfunctor $M_{\mathrm{tor}}$. The quotient $M/M_{\mathrm{tor}}$ is objectwise torsionfree. Both ends of the short exact sequence
\[
0\longrightarrow M_{\mathrm{tor}}\longrightarrow M
\longrightarrow M/M_{\mathrm{tor}}\longrightarrow0
\]
are acyclic by the preceding paragraphs, hence so is $M$. In particular $T_1(M)=0$ for every $M$. Since $\Ab^A$ has enough projectives and colimit is right exact, this is exactness of $\colim_A$. Empty and disconnected categories cause no exception in the argument.
\end{proof}

\begin{proposition}[Proved: the resolving closure and its tensor operation]\label{A:closure}
For a total category $A$, let $\mathcal R_A$ be the smallest full class of objects of $\Ab^A$ containing all $\Z[H]$, $H:A\to\Set$, and closed under extensions, kernels of epimorphisms between its objects, direct sums, direct summands, and filtered colimits. Every object of $\mathcal R_A$ is colimit-acyclic. Moreover, for every $H:A\to\Set$, the functor $M\mapsto M\otimes_\Z\Z[H]$ preserves $\mathcal R_A$. Adding the tensor operation to the definition therefore does not enlarge this class.
\end{proposition}
\begin{proof}
The bar complex of $\Z[H]$ is the integral nerve-chain complex of $\El_AH$, so totality proves the initial vanishing. Long exact sequences give the extension and kernel closures, and the bar complex gives the other three closures. For fixed $H$, pointwise tensoring with $\Z[H]$ is exact, preserves direct sums, retracts, and filtered colimits, and sends $\Z[G]$ to $\Z[G\times H]$. The class of all $M$ for which $M\otimes\Z[H]\in\mathcal R_A$ consequently contains the generators and has every closure used to define $\mathcal R_A$. Minimality gives the assertion. We have not asserted tensor stability for the larger class of all colimit-acyclic diagrams.
\end{proof}

\begin{remark}[Unresolved: the membership step]\label{A:membership}
For an abelian diagram $M$, the augmentation ideal $I(M)=\ker(\Z[UM]\to\underline\Z)$ and all its tensor powers lie in $\mathcal R_A$. The natural epimorphism $I(M)\to M$, $[m]-[0]\mapsto m$, has kernel $I(M)^2$, where multiplication is in the group ring of the additive group $M(a)$ at each object. Although $I(M)^2$ is the image of $I(M)\otimes I(M)\to I(M)$, images are not among the proved closure operations. No acyclicity of that image, or membership of the column modules in $\mathcal R_A$, is inferred here.
\end{remark}
